\chapter{Постановка задачи}
\label{cha:Postanovka}

Основные задачи решаемые в рамках проекта:

 \section{Укладка множества коробок в паллету}
 Необходимо разместить все коробки из множества $S$ без взаимных пересечений, максимизируя коэффициент заполнения и одновременно соблюдая ограничения по устойчивости. В первой версии устойчивость игнорируем. 

 \section{Укладка множества коробок в несколько паллет}
 Часто происходят ситуации, когда груз поступает партиями разных адресов, поэтому необходимо разбить исходный массив на подмножества по адресу доставки и для каждого решить подзадачу укладки, минимизируя дисперсию высот, полученных паллет.

 
 \section{Укладка множества коробок в одну или несколько паллет с учетом параметра Aisle, описывающего группы товаров по секциям в магазинах }
 
При укладке паллеты необходимо ввести поощрения за компактное размещение товаров одной группы внутри паллеты. 

Вспомогательные задачи: 

  \section{Разработка программного решения для анализа массива данных }

Должен быть разработан инструмент анализа датасета. Результатом должен быть набор метрик $M_k$. Далее метрики должны использоваться для тонкой подстройки алгоритмов. 

\section{Разработка синтетического  массива данных }

Должно быть разработано программное решение для создания синтезированного массива данных по заданному набору значений метрик. 


  